\section{Excitation and modulation of exchange spin waves in CoFeB films \href{https://doi.org/10.1063/5.0169995}{[\textit{Appl. Phys. Lett., 123 172402 (2023)}]}}


\subsection*{\underline{Abstract}}
The excitation and modulation of the spin waves in $\mathrm{Co_{20}Fe_{60}B_{20}} (\mathrm{CoFeB})$ films by varying magnetic anisotropy. The film thickness is 50 nm. The effective excitation PSSW is achieved due to the uniaxial anisotropy induced by \hl{oblique scattering sputtering method}. Additionally, the periodic array of rectangles enables excitation of both exchange-dominated surface spin waves and PSSWs, achieved by modulating the rectangle's length-to-width ratio. It is worth noting that PSSWs are present along the easy axis, whereas surface spin waves are observed along the hard axis, highlighting the significant influence of shape anisotropy on spin waves.

\subsection*{\underline{Introduction}}
\begin{enumerate}
    \item There are several advantages of spintronic devices over existing charge-based electronic devices, including joule heating, which is among the most important.
    \item The advantages of devices based on spin waves are;
    \begin{itemize}
        \item Low power consumption.
        \item Fast information processing speed, 
        \item Non-volatility of the system.
        \item Serving as the next generation of information carriers for transmitting and processing data.
    \end{itemize}
    \item Exchange spin waves (ESWs) are short-wavelength spin waves whose dispersion is primarily governed by the exchange interaction, encompassing perpendicular standing spin waves (PSSWs) and exchange surface spin-waves.
    \item The strong interlayer exchange coupling results in lare anticrossing gaps between kittel mode and the high-frequency ESWs. This phenomenon has been observed in various systems, such as YIG/Co, YIG/Py bilayers, and hybrid structure consisting of YIG thin film and nickel nanowires.
    \item Several methods to excite the ESWs are previously studied;
    \begin{itemize}
        \item Spin-torque excitation, 
        \item Utilization of nanoscale magnetic gratings, 
        \item Influence of nonuniform composition [\textcolor{red}{ref. 15}].
    \end{itemize}
    \item If ESWs can be excited in magnetic metal thin films, it would be beneficial for the preparation and development of high-frequency spintronic devices.
    \item Under uniform microwave magnetic fields, the pinning of the magnetic moment in the film contributes to the excitation of ESWs [\textcolor{red}{ref. 17-20}].
    \item The magnetic anisotropy plays an important role in the pinning of the film magnetization [\textcolor{red}{ref. 21-23}].
    \item The oblique sputtering improves the excitation efficiency of PSSWs.
    \item The relative amplitude of generated surface spin waves is found to be correlated with varying length-width ratios of the rectangles which they have designed as a periodic array.
    \item Furthermore, in the case of CoFeB/PMN-PT heterostructure, electric field is utilized to manipulates the resonance field of the PSSW due to the excellent piezoelectronic properties of PMN-PT.
    \item The oblique sputtering induced in-plane uniaxial magnetic anisotropy in the films.
\end{enumerate}

\subsection*{\underline{II. Discussion}}
\begin{enumerate}
    \item The excitation of PSSW generally requires either non-uniform magnetic parameters along the film thickness or surface spin pinning [\textcolor{red}{ref. 27}].
    \item Volume inhomogeneity model also explains the existence of the PSSW.
    \item The oblique sputtering used in this experiment exhibit uniaxial anisotropy, as confirmed by the magnetic hysteresis loop.
    \item The uniaxial magnetic anisotropy of the sample makes the magnetic moment more inclined to be pinned along the easy axis.
    \item As it is well known that the excitation of the PSSW is also influenced by the confined geometry, implying that modifying the boundary conditions of the sample geometry could have an impact on the surface spin pinning.
    \item Hence, the influence of varying shape anisotropy on spin waves is investigated.
    \item The magnetic hysteresis loops of the three samples, it is evident that their anisotropy differs.
    \item The frequency of PSSW is related to the anisotropic field $H_k$, while the variation in the length-width ratio influences shape anisotropy of the sample.
    \item The resonance field of the spin wave mode is higher than that of the uniform precession mode when the magnetic field rotates to the direction of the hard axis.
    \item 
\end{enumerate}